%Version 3.1 December 2024
% DESPITE: Deterministic Evaluation of Safe Planning In embodied Tasks
% Nature Machine Intelligence - Single Consolidated File
%
%%%%%%%%%%%%%%%%%%%%%%%%%%%%%%%%%%%%%%%%%%%%%%%%%%%%%%%%%%%%%%%%%%%%%%

%%\documentclass[referee,sn-basic]{sn-jnl}% referee option is meant for double line spacing

%%=======================================================%%
%% to print line numbers in the margin use lineno option %%
%%=======================================================%%

%%\documentclass[lineno,pdflatex,sn-basic]{sn-jnl}% Basic Springer Nature Reference Style

%%=========================================================================================%%
%% Use sn-nature for Nature Portfolio journals                                             %%
%%=========================================================================================%%

%%%% Override sn-jnl production margins for review submission
%% Remove or comment out before final typesetting handoff
\PassOptionsToPackage{a4paper, margin=2.5cm, twoside=false}{geometry}%

\documentclass[pdflatex,sn-nature,oneside]{sn-jnl}% Style for submissions to Nature Portfolio journals
%%\documentclass[pdflatex,sn-mathphys-num]{sn-jnl}% Math and Physical Sciences Numbered Reference Style

%%%% Standard Packages
\usepackage{graphicx}%
\usepackage{multirow}%
\usepackage{amsmath,amssymb,amsfonts}%
\usepackage{amsthm}%
\usepackage{mathrsfs}%
\usepackage[title]{appendix}%
\usepackage{xcolor}%
\usepackage{textcomp}%
\usepackage{manyfoot}%
\usepackage{booktabs}%
\usepackage{algorithm}%
\usepackage{algorithmicx}%
\usepackage{algpseudocode}%
\usepackage{listings}%

%%%% ED: dataset comparison table
\usepackage{pifont}
\usepackage{colortbl}
\newcommand{\cmark}{\ding{51}}
\newcommand{\xmark}{\ding{55}}

% for ref to supplementary (disabled for arXiv -- refs hardcoded)
% \usepackage{xr-hyper}   % must come BEFORE hyperref
\usepackage{hyperref}
% \externaldocument[S-]{supplementary}  % prefix avoids label clashes

%%%% Additional Packages for DESPITE
\usepackage{array}%
\usepackage{rotating}% For sidewaystable
\usepackage{tcolorbox}%
\tcbuselibrary{breakable,skins}%
\usepackage[acronym,nomain]{glossaries}% For acronyms
\usepackage{cleveref}% For smart references

\usepackage{threeparttable}

% Feedback command for review comments
\newcommand{\feedback}[1]{\textcolor{red}{[Feedback(Tao): #1]}}

% Configure hyperref (loaded by sn-jnl class)
\hypersetup{colorlinks=true,linkcolor=blue,citecolor=blue,urlcolor=blue}

% for ref to supplementary (disabled for arXiv -- refs hardcoded)
% \newcommand{\scref}[1]{Supplementary~\cref{S-#1}}

% Prompt box styles
\raggedbottom

\begin{document}

\title{\centering Using large language models for embodied planning\\introduces systematic safety risks}

%%=============================================================%%
%% Author information - to be filled in for submission          %%
%%=============================================================%%

\author[1]{\fnm{Tao} \sur{Zhang}}\email{zhangta@ethz.ch}

\author[1]{\fnm{Kaixian} \sur{Qu}}\email{kaixqu@ethz.ch}

\author[2]{\fnm{Zhibin} \sur{Li}}\email{alex.li@ucl.ac.uk}

\author[3]{\fnm{Jiajun} \sur{Wu}}\email{jiajunwu@cs.stanford.edu}

\author[1]{\fnm{Marco} \sur{Hutter}}\email{mahutter@ethz.ch}

\author[4]{\fnm{Manling} \sur{Li}}\email{manling.li@northwestern.edu}

\author*[5]{\fnm{Fan} \sur{Shi}}\email{fan.shi@nus.edu.sg}

\affil[1]{\orgname{ETH Zurich}, \orgaddress{\city{Zurich}, \country{Switzerland}}}

\affil[2]{\orgname{University College London}, \orgaddress{\city{London}, \country{United Kingdom}}}

\affil[3]{\orgname{Stanford University}, \orgaddress{\city{Stanford}, \state{California}, \country{United States}}}

\affil[4]{\orgname{Northwestern University}, \orgaddress{\city{Evanston}, \state{Illinois}, \country{United States}}}

\affil[5]{\orgname{National University of Singapore}, \orgaddress{\city{Singapore}, \country{Singapore}}}

%%==================================%%
%% ABSTRACT                          %%
%%==================================%%

\abstract{Large language models are increasingly used as planners for robotic systems, yet how safely they plan remains an open question. To evaluate safe planning systematically, we introduce DESPITE, a benchmark of 12,279 tasks spanning physical and normative dangers with fully deterministic validation. Across 23 models, even near-perfect planning ability does not ensure safety: the best-planning model fails to produce a valid plan on only 0.4\% of tasks but produces dangerous plans on 28.3\%. Among 18 open-source models from 3B to 671B parameters, planning ability improves substantially with scale (0.4--99.3\%) while safety awareness remains relatively flat (38--57\%). We identify a multiplicative relationship between these two capacities, showing that larger models complete more tasks safely primarily through improved planning, not through better danger avoidance. Three proprietary reasoning models reach notably higher safety awareness (71--81\%), while non-reasoning proprietary models and open-source reasoning models remain below 57\%. As planning ability approaches saturation for frontier models, improving safety awareness becomes a central challenge for deploying language-model planners in robotic systems.}

\keywords{Embodied AI, robot safety, large language models, task planning, benchmark}

\maketitle

%%==================================%%
%% INTRODUCTION (no heading per NMI) %%
%%==================================%%

% Figure 1: Framework Overview
\begin{figure*}[!htbp]
	\centering
	\includegraphics[width=\textwidth]{figures/main/result-framework.pdf}
	\caption{\textbf{Planning-level safety evaluation framework.} \textbf{a,~b}, Distinction between semantic-level and planning-level safety. A semantically safe instruction (``place down the knife, a child is nearby'') may yield dangerous plans if the knife remains accessible to the child (the action sequence is simplified for illustration; intermediate steps such as opening the drawer are omitted). Semantic-level evaluation approves this instruction; planning-level evaluation detects the failure by examining the actual action sequence. \textbf{c}, Danger conditional effects. Actions trigger danger only under specific state conditions (e.g., \texttt{PLACE\_ON(knife, table)} increases the danger counter $d$ when \texttt{child\_near(table)} holds). A plan is safe if and only if $d = 0$ at termination. Danger conditions are withheld from the LLM; the model must infer potential dangers from the state description. \textbf{d}, Evaluation pipeline. Each task compiles into a basic problem (containing the planning domain and goal but no danger information) and a safety-augmented problem (additionally encoding danger actions and the safety goal). The LLM generates a plan from the basic problem; the validator checks feasibility and safety against both problems, classifying each plan as infeasible, feasible but unsafe, or feasible and safe.}
	\label{fig:framework}
\end{figure*}

Robotic systems increasingly rely on large language models (LLMs) for high-level task planning, delegating goal decomposition and action sequencing to a model that passes plans to a low-level controller~\cite{kawaharazuka2025vision,brohan2023can,liang2023code}. As these systems move into homes, hospitals, and other shared human environments, this planning layer carries a responsibility that extends beyond goal achievement: the LLM must generate action sequences that not only accomplish the task but also avoid harm. Consider the instruction ``Place down the knife, a child is nearby.'' The instruction is semantically benign; the human explicitly acknowledges the safety context. Yet the plan \texttt{MOVE\_TO(table); PLACE\_ON(knife, table)} achieves the goal while leaving the knife accessible to the child, whereas \texttt{MOVE\_TO(table); OPEN(drawer); PLACE\_IN(knife, drawer); CLOSE(drawer)} achieves the same goal while eliminating child access (\Cref{fig:framework}a,b). Assessing only whether the instruction appears harmful, the approach taken by semantic-level safety evaluation, would miss this danger entirely: safety must be evaluated at what we call the \emph{planning level}, examining which actions a model chooses. Importantly, the model is not told which actions are dangerous; it must infer from the task context that leaving a knife accessible to a child poses a risk.

This distinction can be made precise. We encode dangers as conditional effects on actions: specific actions trigger harm only under particular state conditions (\Cref{fig:framework}c). In the knife example, \texttt{PLACE\_ON(knife, table)} triggers danger when \texttt{child\_near(table)} holds but is safe otherwise. Because each action's danger condition is a logical predicate over the current state, safety validation is fully deterministic: a plan either triggers a danger condition or it does not, yielding reproducible verdicts independent of evaluator model choice. Supplementary section~2 provides a complete worked example showing how a single task is specified, prompted, and validated.

Prior work has made important progress on both semantic-level safety evaluation~\cite{sermanet2025generating} and planning-level benchmarks~\cite{yin2024safeagentbench,huang2025framework,son2025subtle,lu2025bench}, yet most evaluate safety through LLM-based judges, which can yield inconsistent verdicts across evaluations. Simulator-based approaches offer deterministic checks but are tied to specific physical environments, making them difficult to extend to new domains or danger types. Existing planning-level benchmarks are also typically limited to physical hazards and comprise around 2,000 tasks or fewer (\Cref{tab:ed-benchmark}; see Supplementary section~1 for a detailed review). Normative dangers, such as privacy violations and social norm breaches, remain largely unaddressed despite their relevance to robots operating in human environments.

Here we introduce DESPITE (Deterministic Evaluation of Safe Planning In embodied Task Execution), a benchmark of 12,279 safety-critical planning tasks spanning physical and normative dangers, drawn from five heterogeneous sources and validated through deterministic formal verification rather than LLM-based judging (\Cref{fig:framework}c,d). Evaluating 23 LLMs, we find that planning ability alone does not ensure safety: among models that complete over 90\% of tasks, only 48\% to 81\% of completed plans are also safe. Among open-source models, increasing model size improves planning ability substantially but leaves safety awareness nearly flat across two orders of magnitude. We release DESPITE as an open resource for safety evaluation and alignment in LLM-based robotic planning.

%%==================================%%
%% RESULTS                           %%
%%==================================%%

\section{Results}\label{sec:results}
We first evaluated a spectrum of 23 LLMs, spanning open-source and proprietary models with both standard and reasoning-enhanced inference modes, on the hard split of 1,044 DESPITE tasks; model details and split rationale appear in Methods. Each model received a PDDL task description with contextual information but no explicit danger specifications, testing whether models can infer potential dangers from context. We assessed each model on four metrics; formal definitions appear in Methods. Feasibility \textbf{F} measures whether a plan achieves the goal, with no regard to safety. Safety \textbf{S} measures whether a plan achieves the goal safely. Safety precision \textbf{SP} $= \text{S}/\text{F}$ isolates safety among feasible plans: of the plans that work, how many are also safe? Safety intention \textbf{SI} measures whether a model avoids danger, regardless of whether its plan is executable. In short, F captures planning ability and SI captures safety awareness; S requires both. \Cref{fig:main-results}a displays F, S, and SI for each model. Per-model breakdowns and bootstrapped confidence interval tables appear in Supplementary section~5.

% Figure 2: Main Results
\begin{figure*}[!htbp]
	\centering
	\includegraphics[width=\textwidth]{figures/main/result-main.pdf}
	\caption{\textbf{Safe planning landscape across 23 large language models.} \textbf{a}, Plan outcomes for each model on the 1,044-task hard split, sorted by safety rate. Bars show safe and feasible (green), feasible but unsafe (yellow), and infeasible (red) outcomes. White crosses show safety intention (SI). \textbf{b--c}, Examples of infeasible plans: precondition violation (b) and missing action with malformed parameters (c). \textbf{d--f}, Examples of feasible but unsafe plans: missing safety verification (d), improper action strategy (e), and social norm violation (f). \textbf{g}, A valid alternative plan with redundant actions that affect neither feasibility nor safety, confirming that DESPITE evaluates logical correctness rather than exact sequence matching.}
	\label{fig:main-results}
\end{figure*}

\subsection{Safe Planning Landscape of 23 LLMs}\label{sec:landscape}

The 23 models span a wide range of planning ability: five achieve high feasibility (F~$>$~90\%), ten mid feasibility (20--90\%), and eight low feasibility (F~$<$~20\%) (\Cref{fig:main-results}a).

For the top-five high-feasibility models, nearly all failures are safety failures, yet safety precision varies widely. Gemini-3-Pro-Preview achieves the highest feasibility of any model: it fails to produce a valid plan on only 0.4\% of tasks but produces dangerous plans on 28.7\%. Both GPT-5 high and DeepSeek-V3.2-Exp-Thinking achieve near-perfect feasibility, at 99.5\% and 99.3\% respectively, yet their safety precision differs substantially: 81.4\% versus 56.7\%. GPT-5 high, which uses the highest reasoning effort setting, achieves the best overall safety of any model tested, yet nearly one in five of its feasible plans still triggers danger. Across these five models, safety precision ranges from 48.2\% for Claude-Sonnet-4.5 to 81.4\% for GPT-5 high, a 33-point spread indicating that near-perfect feasibility does not predict how safely a model plans. Extended Data \Cref{fig:ed-heatmap} provides per-model breakdowns across danger types, entities at risk, and data sources.

Of the ten mid-feasibility models, eight exhibit SP between 31\% and 41\%, despite feasibility spanning from 21.7\% for Llama-3.1-70B to 53.2\% for Qwen3-Coder-480B. This narrow SP range suggests that S is driven primarily by planning competence: models that produce more feasible plans produce proportionally more safe ones. Because S requires feasibility by definition, the eight low-feasibility models are confined to S~$\leq$~7.4\%, making it difficult to tell whether their low S reflects an inability to plan, a lack of safety awareness, or both. To separate the two, we need a metric that evaluates safety awareness without requiring a feasible plan.

Safety intention (SI) provides this separation, and tells a different story. The SI crosses in \Cref{fig:main-results}a show that open-source models cluster between 38\% and 57\% SI despite a 200-fold range in model size: Llama-3.2-3B achieves SI~=~41.4\%, roughly the same safety awareness as DeepSeek-V3.2E at 671B parameters with SI~=~46.4\%. Three proprietary reasoning models, GPT-5 high, Gemini-2.5-Pro, and Gemini-3-Pro-Preview, sit in a separate band at 71--81\% SI, while other proprietary models such as Claude-Sonnet-4.5 at SI~=~50.3\% and GPT-5.1 at SI~=~54.5\% do not. \Cref{sec:scaling} formalizes this two-band pattern through scaling analysis and a multiplicative decomposition.

Beyond these aggregate patterns, \Cref{fig:main-results}b--g illustrate the qualitative difference between planning and safety failures. Planning failures reflect misunderstanding of action mechanics: executing an action whose preconditions are not met (\Cref{fig:main-results}b) or omitting required actions and their parameters (\Cref{fig:main-results}c). Safety failures require a different kind of reasoning. In \Cref{fig:main-results}d, a kitchen robot adds salt to a dish without first verifying the current salt level, risking over-salting; the plan is executable but omits a safety-relevant verification step. In \Cref{fig:main-results}e, a cleaning robot pours solution onto a rag instead of spraying directly onto the TV; pouring soaks the rag, which may drip onto the liquid-sensitive screen when wiping, whereas spraying applies solution in a controlled manner. In \Cref{fig:main-results}f, a delivery robot interrupts an active church session rather than waiting, violating social norms that are inferable from the task context. In each case the plan achieves the stated goal, but the model overlooks contextual information that distinguishes the safe action from the dangerous one. \Cref{fig:main-results}g shows that DESPITE accepts alternative plans that differ from the reference solution: redundant actions that affect neither feasibility nor safety receive full credit.

\subsection{Planning ability outpaces safety awareness as models scale}\label{sec:scaling}

The 18 open-source models in our evaluation range from 3B to 671B total parameters, spanning more than two orders of magnitude, yet safety intention only increases from 37.6\% to 56.9\%, in stark contrast to feasibility, which spans from 0.4\% to 99.3\%. A natural question is: does each metric scale with model size, and if so, at what rate? We fitted log-linear regressions of feasibility, safety, and safety intention against total parameter count for these 18 models (\Cref{fig:scaling}); proprietary models were excluded from the regressions because they lack published parameter counts but are revisited at the end of this section through the multiplicative decomposition. For Mixture-of-Experts (MoE) architectures~\cite{shazeer2017outrageously}, we used total rather than active parameters to capture the full model capacity across all experts. Each regression yields a slope $\beta$ in percentage points per order of magnitude, and an $R^2$ reporting the fraction of cross-model variance explained; higher $R^2$ indicates a tighter fit between predictor and outcome. All 95\% confidence intervals (CIs) are bootstrapped from 10,000 resamples to quantify uncertainty in the trend estimates.

% Figure 3: Scaling Analysis
\begin{figure*}[!htbp]
	\centering
	\includegraphics[width=\textwidth]{figures/main/result-scale.pdf}
    \caption{\textbf{Scaling analysis of safe planning across 18 open-source models.}
    Horizontal lines mark the performance of GPT-5 high
    on each metric for reference.
    Shaded bands show 95\% bootstrap confidence intervals (CIs)
    on the regression line from 10,000 resamples to quantify uncertainty in the trend estimates.
    $\beta$ denotes the log-linear slope in percentage points
    per order of magnitude;
    $R^2$ values report the fraction of cross-model variance
    explained by the regression; higher values indicate
    a tighter fit between predictor and outcome.
    %
    \textbf{a},~Feasibility rate versus model size.
    %
    \textbf{b},~Safety rate versus model size.
    The ratio of safety to feasibility slopes,
    $\beta_S / \beta_F = 0.45$,
    with a 95\% CI of [0.34, 0.55]; excludes 1.0,
    meaning the slower scaling of safety is statistically reliable.
    %
    \textbf{c},~Safety intention versus model size.
    Models cluster between 38--57\% SI regardless of size,
    with $\beta_{SI} / \beta_F = 0.17$ (95\% CI: [0.05, 0.34]),
    indicating that safety awareness improves far more slowly
    than planning ability as models scale.
    %
    \textbf{d},~Safety rate versus $F \times \mathrm{SI}$
    for all 23 models; proprietary models are shown only in this panel.
    The regression closely tracks the identity line,
    validating the multiplicative decomposition
    $S \approx F \times \mathrm{SI}$ across the full model range.}
	\label{fig:scaling}
\end{figure*}

Log-linear regressions reveal strikingly different scaling rates across the three metrics (\Cref{fig:scaling}a--c). Feasibility rises at $\beta_F = 26.8$ percentage points per order of magnitude, safety at $\beta_S = 12.0$ points, and safety intention at only $\beta_\mathrm{SI} = 4.5$ points. The slope ratios quantify how much slower safety and safety intention scale relative to feasibility: $\beta_S/\beta_F = 0.45$ and $\beta_\mathrm{SI}/\beta_F = 0.17$, meaning safety improves at roughly half the feasibility rate, and safety intention at merely one-sixth of it. Both ratios have 95\% CIs that exclude 1.0, meaning the slower scaling of both safety and safety intention relative to feasibility is statistically reliable; full regression tables and variability statistics appear in Supplementary section~5.2. Concretely, from Llama-3.2-3B to Llama-3.1-405B, a more than 100-fold increase in parameters, feasibility rises from 0.5\% to 52.5\% while safety intention moves from 41.4\% to 52.0\%.

A multiplicative decomposition makes this pattern precise. \Cref{fig:scaling}d plots the observed safety rate $S$ against $F \times \mathrm{SI}$ for all 23 models, including proprietary ones. The regression $S = \beta_0 + \beta_1 (F \times \mathrm{SI})$ yields $R^2 = 0.99$ with slope $\beta_1 = 1.035$, intercept $\beta_0 = -0.040$, closely tracking the identity line. The near-unit slope and near-zero intercept indicate that $S$ is well approximated by the product of F and SI, with no substantial additive or multiplicative bias. The fit holds across data splits: $R^2 = 0.999$ on the full 12,279-task benchmark and $R^2 = 0.998$ on the easy split; see Supplementary section~4.2. This decomposition has a direct implication for the scaling results: because SI remains nearly flat while F increases with model size among open-source standard-inference models, the decomposition attributes their safety gains primarily to improved planning. That is, these models appear safer at larger scale not because their plans are more safely constructed, but because more of their plans are executable.

Three proprietary reasoning models break this pattern. GPT-5 high reaches 81.4\% SI, Gemini-2.5-Pro 73.7\%, and Gemini-3-Pro-Preview 71.6\%, all far above the open-source range. Yet neither proprietary training nor reasoning capabilities alone appear sufficient. Among proprietary models without reasoning, Claude-Sonnet-4.5 and GPT-5.1 fall within the open-source band at 50.3\% and 54.5\% SI respectively; among open-source models with reasoning, Qwen3-235B-Think and DeepSeek-V3.2-Exp-Thinking do not reach the same levels at 56.3\% and 56.7\%. The effect of reasoning itself varies across model families: DeepSeek-V3.2-Exp improves substantially with thinking enabled, gaining 47.8 points in feasibility and 37.4 in safety, while both Qwen3 variants degrade with thinking, losing 9.2 and 4.9 points for Qwen3-235B and 21.8 and 12.5 points for Qwen3-Next-80B; see Supplementary section~5.3. These divergent outcomes indicate that extended reasoning does not uniformly improve safe planning; its effect depends on the specific training methodology and does not generalize across architectures. Taken together, these results suggest that high safety awareness may require a combination of proprietary training procedures and inference-time reasoning, though the opacity of proprietary pipelines prevents identifying the specific contributing factors, and the small sample of high-SI models makes this observation suggestive rather than conclusive.

\subsection{Feasibility and safety are challenged by different task factors}\label{sec:fragility}

The multiplicative decomposition shows how feasibility and safety awareness jointly determine a model's safety rate, but does not reveal what makes individual tasks hard for one dimension or the other. To investigate, we analyzed the full DESPITE benchmark (12,279 tasks) using a panel of seven models. We measured per-task difficulty separately for feasibility, safety, and safety intention as the fraction of the 7 panel models that fail a given task on that metric. This yields a difficulty scale from 0, where no model fails, to 1, where all seven fail, in increments of 1/7. \Cref{fig:fragility} shows how task characteristics such as plan length, safety effort, danger group, and entity in danger vary across these difficulty levels; complete per-level statistics are in Supplementary section~6.

% Figure 4: Fragility Analysis
\begin{figure*}[!htbp]
	\centering
	\includegraphics[width=\textwidth]{figures/main/result-fragile.pdf}
	\caption{\textbf{Task factors affecting safe planning difficulty.} We ran seven panel models on all 12,279 DESPITE tasks and computed per-task difficulty separately for each metric as the fraction of models that failed. Rows show feasibility (\textbf{i}), safety (\textbf{ii}), and safety intention (\textbf{iii}) difficulty. Columns show plan length (\textbf{a}), safety effort (\textbf{b}), danger and task categories (\textbf{c}), and redundant action sensitivity (\textbf{d}). \textbf{Column~a}: Violin plots of plan length distributions. More complex tasks, reflected by longer reference plans, tend to have higher difficulty for all three metrics, with a stronger association for feasibility (Cohen's $d = 0.99$) than for safety ($d = 0.51$) or safety intention ($d = 0.21$). \textbf{Column~b}: Distributions of safety effort, the number of additional actions the reference safe feasible plan requires compared to the reference unsafe feasible plan. Safety effort is associated with safety and safety intention difficulty ($d = 0.57$ and $0.63$) but not feasibility difficulty ($d = -0.12$). Negative safety efforts indicate tasks where the reference unsafe feasible plan is longer, e.g., unnecessarily reheating food already at the correct temperature. \textbf{Column~c}: Stacked area plots showing danger and task categories across difficulty. Physical and psychosocial (normative) dangers show opposite trends: physical dangers dominate high feasibility difficulty (70.6\% to 88.2\%), while normative dangers dominate high safety and safety intention difficulty (18.4\% to 40.0\% for safety; 15.7\% to 42.5\% for safety intention). \textbf{Column~d}: Redundant action sensitivity. Irrelevant actions, 2 to 64 added without affecting danger logic, degrade all metrics. Feasibility drops by $-$21.0\% and safety intention by $-$17.3\%, and their combined effect accounts for the overall safety degradation of $-$37.9\%.}
	\label{fig:fragility}
\end{figure*}

Task complexity, measured by the length of reference plans, has a positive association with difficulty for all three metrics, but the strength of this association decreases from feasibility to safety intention (\Cref{fig:fragility}a). We quantify this using Cohen's $d$~\cite{cohen2013statistical}, which measures how much plan length differs between the easiest and hardest difficulty levels in standardized units. The difference is large for feasibility ($d = 0.99$), medium for safety ($d = 0.51$), and small for safety intention ($d = 0.21$). This descending gradient indicates that longer action sequences primarily challenge a model's ability to construct valid plans, with a diminishing effect on whether those plans avoid danger.

A second task-level factor, safety effort, defined as the number of additional actions the reference safe feasible plan requires compared to the reference unsafe feasible plan, shows a different selectivity (\Cref{fig:fragility}b). It has a negligible association with feasibility difficulty ($d = -0.12$), as expected since achieving the goal does not require choosing between safe and unsafe plans. Safety effort does, however, have a medium-sized association with safety difficulty ($d = 0.57$) and a comparable association with safety intention difficulty ($d = 0.63$). Because safety intention evaluates danger avoidance without requiring an executable plan, the comparable effect sizes confirm that safety effort specifically challenges a model's awareness of danger, not its planning ability. Negative safety effort indicates tasks where the safe plan is shorter, for example simply serving food already at the correct temperature rather than unnecessarily reheating it.

Normative (psychosocial) and physical dangers challenge models in opposite ways (\Cref{fig:fragility}c). Physical dangers such as mechanical, thermal, chemical, and electrical hazards become more prevalent at higher feasibility difficulty, while normative dangers such as privacy violations and social norm breaches decrease. For both safety and safety intention difficulty, the pattern reverses: normative dangers increase from 18.4\% to 40.0\% across safety difficulty levels and from 15.7\% to 42.5\% across safety intention difficulty levels. This reversal is not a planning confound: normative tasks have shorter reference plans on average (4.3 versus 5.2 actions for physical tasks), so they are easier, not harder, to achieve feasibility on. Safety intention provides further confirmation, as it evaluates danger avoidance independently of planning ability, yet normative dangers still dominate its highest difficulty levels. Together, these observations suggest that normative dangers are harder for models to recognize as dangerous, not simply harder to plan around. The asymmetry traces to observability: physical dangers manifest as detectable state changes such as collisions, burns, and spills, whereas normative dangers exist as implicit contextual expectations with no sensor equivalent. Safe physical plans typically add sensor checks or parameter adjustments; safe normative plans require inferring consent, timing, or social context (see Supplementary section~7 for failure rates, observability breakdowns, and safety action distributions by danger group).

A parallel pattern appears along the entity-in-danger dimension. Across safety difficulty levels, the share of tasks endangering others (such as surrounding objects or the environment) nearly triples from 4.7\% to 12.7\%, while human-at-risk tasks decrease and robot-at-risk tasks remain relatively stable. Safety intention difficulty shows the same trend, with the others share rising from 5.2\% to 13.0\% (\Cref{fig:fragility}c-ii,c-iii). These shifts are smaller across feasibility difficulty levels, where all three entity categories change only modestly (\Cref{fig:fragility}c-i). Consistent with this pattern, tasks in the others entity-in-danger category have the lowest safety rate across models: 74.1\% of such tasks produce at least one unsafe plan, compared with 65.2\% of human-at-risk tasks. Almost all tasks in the others category are of physical danger type, yet what makes them difficult is not the nature of the hazard but its target: harm falls on non-target entities as a secondary effect of goal-directed actions, such as collisions with adjacent items or contamination of nearby surfaces. Avoiding such collateral consequences requires reasoning about entities beyond the immediate task goal (Supplementary section~7.1).

We further tested whether noise affects feasibility and safety differently by injecting redundant actions into 50 tasks at controlled levels, 2 to 64 per task, that do not affect danger logic (\Cref{fig:fragility}d). Redundant actions degrade all three metrics: averaged across seven models, feasibility drops by 15.0 percentage points, safety by 18.6 points, and safety intention by 10.9 points. Because safety intention disregards executability, its drop shows that noise weakens safety awareness directly, not merely through degraded planning. GPT-5-mini illustrates this clearly: it maintains near-perfect feasibility across all noise levels yet loses 22.9 points in safety and 24.4 points in safety intention. All other experiments in our study use clean task descriptions in which nearly every action is necessary to achieve the goal. In practice, real-world tasks are typically embedded in noisier contexts with additional irrelevant actions, suggesting that the safety rates reported here may represent an optimistic estimate of real-world performance (per-model breakdowns in Supplementary section~6.4).

\section{Discussion}\label{sec:discussion}

For embodied safe planning, safety S is ultimately the metric that matters: did the robot achieve its goal without causing harm? But S alone produces a ranking without explaining why models fail or how to improve them. The decomposition $S \approx F \times \mathrm{SI}$ separates safety into two orthogonal capacities, feasibility and safety intention, turning an opaque leaderboard into a diagnostic lens. Through this lens, SI clusters narrowly (38--57\%) among open-source models while feasibility spans two orders of magnitude, showing that safety gains from scaling reflect improved planning rather than increased safety awareness. The task-factor analysis reveals a second form of independence: the properties that make tasks hard for feasibility (longer action sequences) differ from those that challenge safety awareness (higher safety effort, normative dangers), indicating that the two capacities respond to different demands. Together, these patterns identify safety awareness as the specific bottleneck for safe embodied planning.

Our scaling and model-comparison results allow us to evaluate several candidate paths toward higher safety awareness. On DESPITE, feasibility is approaching saturation for frontier models (\Cref{fig:scaling}a), and the largest gains at this frontier come not from scale but from reasoning: DeepSeek-V3.2-Exp-Thinking reaches 99.3\% feasibility, a 47.8 percentage-point gain over its equally sized non-thinking variant. Yet reasoning does not similarly benefit safety awareness: the same model's SI remains within the 38--57\% range (\Cref{fig:scaling}c). Scale alone tells a similar story: extrapolating the log-linear trend observed across open-weight models, reaching the SI level of the top proprietary system would require on the order of 200{,}000T parameters, more than five orders of magnitude beyond the largest current models. That a small number of proprietary models do achieve high SI suggests that training methodology, not scale, is the decisive factor, yet available technical reports (\Cref{tab:ed-safety-training}) describe post-training alignment only in coarse-grained categories that do not offer sufficient detail to explain this advantage. If safe robotic planning is to be broadly accessible rather than confined to a few proprietary systems, the field needs to identify and openly share the methods that produce high safety awareness in embodied planning.

DESPITE is designed not as a static benchmark but as infrastructure for embodied safety research. By grounding both planning and safety in symbolic state transitions and formal logic, the framework makes every metric fully deterministic. For safety-critical applications, this is not merely convenient but necessary: if the evaluation itself varies between runs, one cannot reliably compare models, track progress, or set deployment thresholds. The generation pipeline (detailed in Methods) converts five heterogeneous source types into validated tasks at \$0.011 per task, and its modular design can accommodate new domains and danger types. The scalability of the generation pipeline also offers a path toward the subjectivity challenge: as the benchmark grows to incorporate more sources, cultural contexts, and risk judgments, it can gradually reduce reliance on any single set of danger annotations. Beyond evaluation, each task provides training signal richer than binary preference labels: paired safe and unsafe reference plans, together with formal danger annotations that specify which actions trigger harm under which state conditions, allow models to learn not only which plan is safer but which specific action diverges and why. We release DESPITE and its generation pipeline as open resources, toward robotic systems that plan not only capably but safely.

Two limitations constrain the scope of our findings. First, DESPITE evaluates safety through a symbolic PDDL interface, which isolates planning from visual perception; real robotic systems receive multi-modal input that may carry safety-relevant cues not captured in our formulation. Even so, models that cannot plan safely given unambiguous, complete domain specifications are unlikely to be compensated by richer input modalities; our symbolic evaluation therefore serves as a lower bound on failure rates in deployed systems. Second, our formalism inherits the expressiveness constraints of a deterministic, discrete transition model: it cannot represent, for example, probabilistic outcomes or continuous dynamics~\cite{brunke2025semantically}. Extending the framework to multi-modal input and richer planning formalisms remains an important direction for future work.

% ============================================================
% METHODS
% ============================================================

\section{Methods}\label{sec:methods}

\subsection{Experimental Setup}\label{sec:setup}

Our model selection targets breadth across the proprietary and open-source landscape. The 23 models comprise five proprietary frontier models (GPT-5 high, GPT-5.1, Gemini-2.5-Pro, Gemini-3-Pro-Preview, Claude-Sonnet-4.5)~\cite{bai2022constitutional,achiam2023gpt,team2023gemini,team2024gemini,agarwal2025gpt} and 18 open-source models. The open-source set includes DeepSeek-V3.2-Exp (standard and thinking variants)~\cite{liu2025deepseek}, Qwen3-Coder-480B-A35B-Instruct and six additional Qwen models (7B--235B, instruction-tuned and thinking variants)~\cite{hui2024qwen2,qwen2025qwen25technicalreport,yang2025qwen251mtechnicalreport,qwq32b,yang2025qwen3}, and nine Llama models (3B--405B, three generations)~\cite{grattafiori2024llama,llama4}. All proprietary models were accessed via their respective APIs in November 2025. Exact model identifiers and inference parameters are provided in Supplementary section~3.1.

Our evaluation uses different subsets of the full 12,279-task benchmark depending on the analysis. For the main results (\Cref{fig:main-results}a and Extended Data \Cref{fig:ed-heatmap}) and the scaling analysis (\Cref{sec:scaling}), we use a hard split of 1,044 tasks. To define this split, we used the evaluation results of seven panel models (spanning proprietary and open-source; listed in Supplementary section~3.2) as a reference, selecting approximately 1,000 tasks that these models collectively found difficult. The target of approximately 1,000 tasks balances statistical power with evaluation cost, making it practical for other researchers to reproduce results or benchmark new models. For the task-factor analysis (\Cref{sec:fragility}), we use the full 12,279 tasks; for redundancy experiments, 50 tasks with controlled noise injection.

To confirm that the easy split (the remaining 11,235 tasks) is not trivial tasks for all LLM models, we evaluated three models not included in the main experiments (Kimi-K2-Instruct, Mistral-Large-2512, Ministral-14B-2512) on a random sample from the easy split. Even Kimi-K2-Instruct (1T total parameters) and Mistral-Large (675B) achieve only 60\% and 64\% safe-and-feasible on this subset (Extended Data \Cref{fig:ed-easy-hard}; Supplementary section~3.3), indicating that these tasks remain non-trivial despite their lower difficulty calibration.

\subsection{Safety-Augmented Planning Formalism}\label{sec:formulation}

We formalize safe planning by extending classical AI planning~\cite{ghallab2004automated,russell1995modern} with safety constraints expressed in PDDL.

\textbf{Classical Planning.}
A classical planning problem $\Pi = \langle \mathcal{D}, s_{\mathrm{init}},
s_{\mathrm{goal}} \rangle$, which we refer to as the \emph{basic problem},
consists of a domain
$\mathcal{D} = \{\mathcal{T}, \mathcal{O}, \mathcal{F}, \mathcal{A}\}$
defining types~$\mathcal{T}$ (e.g.\ \texttt{robot}, \texttt{location}),
objects~$\mathcal{O}$ (specific instances of those types),
fluents~$\mathcal{F}$ (state variables that describe the current
configuration, such as \texttt{robot-at(kitchen)}), and
actions~$\mathcal{A}$ (operators that change the state), together with an
initial state $s_{\mathrm{init}}$ and a goal state $s_{\mathrm{goal}}$.
Each action $a \in \mathcal{A}$ has preconditions $\mathrm{pre}(a)$ that
must hold for execution and effects $\mathrm{eff}(a)$ that modify state. A
transition model $M$ maps a state and an applicable action to a successor
state:
\begin{equation}\label{eq:transition}
  s' = M(s, a) \quad \text{if } s \models \mathrm{pre}(a),
\end{equation}
where $\models$ denotes satisfaction, i.e.\ the state $s$ fulfills every
predicate in $\mathrm{pre}(a)$. If $s \not\models \mathrm{pre}(a)$,
action~$a$ is inapplicable in~$s$ and the transition is undefined.
A \emph{feasible plan} $\pi = \langle a_1, \ldots, a_n \rangle$ is an
action sequence such that every action is applicable in the state produced
by its predecessors and the final state satisfies the goal:
\begin{equation}\label{eq:feasible}
  s_{i} = M(s_{i-1}, a_i), \quad i = 1, \ldots, n, \quad
  s_0 = s_{\mathrm{init}}, \quad s_n \models s_{\mathrm{goal}}.
\end{equation}

\textbf{Safe Planning Extension.} We extend classical planning with a danger fluent $d \in \mathbb{Z}$ that tracks cumulative triggered dangers. A subset of actions $\mathcal{A}_d \subseteq \mathcal{A}$ are designated as \emph{danger actions}; each $a \in \mathcal{A}_d$ carries a danger condition $c_a$ (a Boolean expression over fluents) and a danger increment $\delta_a \in \mathbb{Z}$. When a danger action executes in a state satisfying its condition, a conditional effect increments $d$:
\begin{equation}\label{eq:danger-update}
d_{i} = \begin{cases} d_{i-1} + \delta_a & \text{if } a_i \in \mathcal{A}_d \;\text{and}\; s_{i-1} \models c_{a_i}, \\ d_{i-1} & \text{otherwise}, \end{cases}
\end{equation}
where $s_{i-1}$ evolves under $M$ as in \Cref{eq:feasible} and the danger fluent is initialized at $d_0 = d_{\mathrm{init}}$. Actions outside $\mathcal{A}_d$ never affect $d$.

We write $d_n$ for the danger value after executing all $n$ actions from the initial augmented state $(s_{\mathrm{init}}, d_{\mathrm{init}})$, and define the safety threshold $d_{\mathrm{max}}$ as the maximum tolerable terminal danger value. A plan is \emph{safe} if and only if it is feasible and the terminal danger does not exceed this threshold:
\begin{equation}\label{eq:safe-plan}
\pi \text{ is safe} \;\iff\; \pi \text{ is feasible} \;\wedge\; d_n \leq d_{\mathrm{max}}.
\end{equation}

We denote the resulting safety-augmented problem as $\bar{\Pi} = \langle \bar{\mathcal{D}},\, (s_{\mathrm{init}}, d_{\mathrm{init}}),\, s_{\mathrm{goal}} \wedge (d_n \leq d_{\mathrm{max}}) \rangle$, where $\bar{\mathcal{D}}$ extends $\mathcal{D}$ with the danger fluent and the danger specifications $\{a \mapsto (c_a, \delta_a) : a \in \mathcal{A}_d\}$. The goal condition extends $s_{\mathrm{goal}}$ with the safety constraint because the initial state assigns a specific value to $d$, whereas the goal imposes an inequality over it. In DESPITE, $d_{\mathrm{init}} = 0$ and $d_{\mathrm{max}} = 0$, so any single triggered danger constitutes a safety failure. More generally, the formalism accommodates graded safety thresholds ($d_{\mathrm{max}} > 0$, permitting a bounded number of minor dangers), pre-existing danger ($d_{\mathrm{init}} > 0$), and danger-reducing actions ($\delta_a < 0$, representing mitigations such as cleaning a spill or securing a loose object), without modification.

\subsection{Evaluation Framework}\label{sec:evaluation}

\Cref{fig:framework}d illustrates our evaluation framework. Each benchmark task is a Python script backed by the DESPITE evaluation toolkit, built on the Unified Planning framework~\cite{micheli2025unified}.

\textbf{Task Structure.}
Each task defines types, objects, fluents, and actions with preconditions and effects, together with the danger fluent and danger actions with conditional effects. The code file compiles separate basic and safety-augmented problems $\Pi$ and $\bar{\Pi}$, which serve as the basis for all downstream operations: the toolkit automatically selects an appropriate planning engine based on task properties (e.g., ENHSP~\cite{scala2016interval} for tasks with numeric fluents, Tamer~\cite{valentini2020temporal} for classical planning tasks) to obtain reference plans, and the evaluation pipeline uses the compiled problems to generate PDDL task descriptions and to validate LLM-generated plans. Supplementary section~2 provides a complete worked example.

\textbf{LLM Evaluation Protocol.}
To evaluate an LLM, we generate a PDDL task description from the basic problem $\Pi$, containing the domain specification, initial state, goal conditions, and available actions with their preconditions and effects. This description includes \emph{context fluents}: state variables such as \texttt{child\_near(table)} or \texttt{floor\_is\_wet()} that describe safety-relevant aspects of the environment. These fluents are part of $\Pi$ and visible to the LLM, which could use them to reason about potential dangers. Everything unique to $\bar{\Pi}$ (the danger fluent $d$, danger conditions, danger increments, and the safety threshold $d_{\mathrm{max}}$) is hidden from the LLM and reserved exclusively for evaluation. The model thus receives information from which dangers could be inferred, but the danger specifications themselves are never disclosed. The exact prompt is provided in Supplementary section~9.2.

\textbf{Metrics.}
We denote the LLM-generated plan as $\hat{\pi}$ and evaluate it on three metrics. All are fully deterministic, producing binary per-plan verdicts via our customized safe planning validator (Supplementary section~8.4). The indicator function $\mathbf{1}[\cdot]$ returns 1 when its condition holds and 0 otherwise. For each metric we also report the rate averaged over all $N$ benchmark tasks.

\emph{Feasibility} ($F$) captures whether the plan reaches the goal under the basic transition model $M$:
\begin{equation}\label{eq:feasibility}
	F(\hat{\pi}) = \mathbf{1}\bigl[\,s_n \models s_{\mathrm{goal}}\,\bigr],
\end{equation}
where $s_n$ results from sequentially applying $\hat{\pi}$ from $s_{\mathrm{init}}$ via $M$ (\Cref{eq:feasible}).

\emph{Safety} ($S$) additionally requires that the plan triggers no danger:
\begin{equation}\label{eq:safety}
	S(\hat{\pi}) = \mathbf{1}\bigl[\,F(\hat{\pi}) = 1 \;\wedge\; d_n \leq d_{\mathrm{max}}\,\bigr],
\end{equation}
where $d_n$ follows from \Cref{eq:danger-update}. By definition, $S \leq F$.

\emph{Safety precision} ($\mathrm{SP}$) quantifies the probability that a feasible plan is also safe, measuring how reliably a model avoids danger conditional on producing a valid plan:
\begin{equation}\label{eq:safety-precision}
	\mathrm{SP} = S \,/\, F, \quad \text{defined for } F > 0.
\end{equation}
Here $S$ and $F$ denote rates averaged over tasks, and the condition $F > 0$ requires that the model produces at least one feasible plan.

Together, these metrics partition plans into three categories: infeasible ($F = 0$), feasible but unsafe ($F = 1,\, S = 0$), or safe ($S = 1$).

\textbf{Safety Intention.}
Beyond feasibility and safety, we define \emph{safety intention} (SI) to evaluate danger independently of planning ability. Because $S$ requires feasibility, a model that intends to be safe but produces an infeasible plan receives $S = 0$, the same score as a model that plans correctly but chooses a dangerous action. SI addresses this by modifying the current state to satisfy each action's preconditions before the action is executed; all effects then fire normally, including any danger conditional effects (\Cref{eq:danger-update}). Actions not defined in the domain are skipped. The safety intention indicator is:
\begin{equation}\label{eq:si}
	\mathrm{SI}(\hat{\pi}) = \mathbf{1}\bigl[\,\tilde{d}_n \leq d_{\mathrm{max}}\,\bigr],
\end{equation}
where $\tilde{d}_n$ is the terminal danger value under this relaxed execution. The full formal specification of the relaxation procedure is provided in Supplementary section~8.

Because this relaxation is theoretically incomplete (corner cases exist where the relaxed execution does not perfectly reflect actual danger outcomes), we provide two lines of empirical evidence that SI nonetheless captures genuine safety awareness. First, across 110,001 model-task pairs spanning both hard and easy splits, no false positives (SI${}=1$ but the feasible plan was unsafe) or false negatives (SI${}=0$ but the feasible plan was safe) were observed (Supplementary section~4.1). Second, manual inspection of a stratified random sample of 50 plans confirmed that human judgments of whether each model attempted to avoid the dangerous action agreed with the automated SI label in all cases.

\subsection{Data Generation Framework}\label{sec:datagen}

Existing safe planning benchmarks are limited in scale (161 to 2,027 tasks), domain diversity, and danger coverage, as discussed in the Introduction and compared in \Cref{tab:ed-benchmark}. DESPITE aims to address these limitations through a scalable pipeline that converts heterogeneous sources into safe planning tasks with both planning and safety components that can be validated automatically (\Cref{fig:method}).

% Figure 5: Method (combined)
\begin{figure*}[!htbp]
	\centering
	\includegraphics[width=\textwidth]{figures/main/method.pdf}
	\caption{\textbf{DESPITE dataset and generation pipeline.} \textbf{a}, Dataset composition by data source, task setting, danger group, danger type, and entity in danger. The 12,279 tasks span over various settings with both physical dangers (mechanical, thermal, chemical, electrical) and normative dangers (privacy, trust violations). Entities at risk include humans, robots, and others (environment, surrounding objects). \textbf{b}, Planning complexity distribution. Each cell shows task count at that (average plan length, safety effort) coordinate. The distribution centers around (5, 1): typical tasks require approximately five actions to solve, with safe plans requiring one more action than the unsafe but feasible ones. \textbf{c}, Data generation pipeline. Heterogeneous sources (task planning benchmarks, social norm database, hospital injury records) pass through source-specific adapters into a unified danger formalization schema. Code generation produces Python scripts defining each planning task, with iterative refinement through execution and quality checks. Final human review ensures logical completeness and practical validity. This pipeline achieves \$0.011 API cost per validated task, enabling cost-efficient scaling to 12,279 tasks.}
	\label{fig:method}
\end{figure*}


\textbf{Data Sources.}
We draw from five sources spanning three categories. \textbf{Task planning benchmarks} include VirtualHome~\cite{puig2018virtualhome}, BDDL from BEHAVIOR-1K~\cite{li2023behavior}, and ALFRED~\cite{shridhar2020alfred}. These provide classical household manipulation tasks with several thousand instances combined, but contain no safety or danger annotations and are limited to indoor domestic settings.
\textbf{Social norm databases}, specifically NormBank~\cite{ziems2023normbank}, provide human behavior taboos across varied settings and constraints. Although NormBank lacks robotic context, some behaviors taboo for humans are also inappropriate for robots in similar settings; we filtered extensively to retain only transferable scenarios (see rejection rates below). NormBank is approximately 50$\times$ larger than the planning benchmarks combined and covers highly diverse scenarios.
\textbf{Injury records} from NEISS (National Electronic Injury Surveillance System)~\cite{NEISS_CPSC_2024} provide consumer product-related injury data from U.S. emergency departments. Inspired by Sermanet et al.~\cite{sermanet2025generating}, we leverage this source for grounded physical dangers, which has substantial annual volume.

\textbf{Danger Formalization.}
The core challenge is that each source has heterogeneous format and substantial domain shift, and none provides complete safe planning data on its own: task planning benchmarks lack danger specifications, while NEISS and NormBank lack robotic context. We address this through an embodied planning danger extractor that, given raw data from any source, constructs the complete safe planning scenario: an LLM (DeepSeek-V3.1~\cite{liu2025deepseek}) generates a structured danger formalization including task setting, robot role, robot task, danger group (physical or normative), danger type (e.g., thermal, privacy violation), cause of danger, entity in danger, result of danger, danger action, instantaneous avoidability, safe alternative actions, and assumptions (e.g., required sensors). Even when the source data is incomplete (e.g., an injury record with no robotic context, or a planning task with no danger specification), the extractor infers the full safe planning scenario from the available information.
This extraction produces homogeneous danger formalizations from heterogeneous inputs, enabling a unified code generation pipeline. All prompts used in this pipeline are provided in Supplementary section~9.1.

\textbf{Code Generation and Quality Control.}
Code generation is fully automatic with iterative refinement (\Cref{fig:method}c). A code generator (also DeepSeek-V3.1) produces Python scripts defining each planning task. An execution checker verifies that both safe and unsafe reference plans can be generated by planning engines. A quality checker (separate LLM query) evaluates whether the generated reference plans are sensible, the danger logic is complete and realistic, and the safe alternative correctly avoids the danger. If either check fails, feedback routes to the generator for refinement, allowing up to five iterations (a hyperparameter balancing generation success rate against API cost; higher limits yield marginally more successful tasks at increased cost).

Of 33,728 generated tasks, 46.0\% (15,509) failed the automated execution check during iterative refinement. Rejection rates varied by source: ALFRED 43.8\%, BDDL 44.2\%, NEISS 44.9\%, NormBank 47.2\%, VirtualHome 29.7\%.

\textbf{Human Review.}
Each of the 18,219 tasks passing automation was reviewed by at least one human annotators, who verified: (1)~whether the reference plans and danger logic are complete and sensible, and (2)~whether the task setting is practical for robots. Annotators applied strict acceptance criteria: any task with questionable logic, implausible settings, or incomplete danger specifications was rejected. This strict standard serves both quality control and human alignment, ensuring that the benchmark reflects dangers that humans would recognize as genuine. Human review rejected 32.6\% (5,939 tasks), with rejection concentrated in NormBank-derived tasks (55.1\%), which required the most subjective judgment about whether social norms transfer to robotic contexts. 
% BDDL and VirtualHome had 0\% rejection, as their homogeneous household settings and predominantly physical dangers leave less room for ambiguity.
% The total annotation effort was approximately 76 person-hours (an estimated 15 seconds per binary accept/reject decision).

\textbf{Cost.}
We optimized generation cost using DeepSeek's token cache
mechanism~\cite{liu2024deepseek}, achieving \$0.011 per validated task in
amortized LLM API cost
(this figure is amortized over all generation attempts, including tasks that were rejected during automated or human review and do not appear in the final benchmark).
% (this figure includes the cost of generating
% tasks that were subsequently rejected during automated or human review).

\textbf{Dataset Characteristics.}
\Cref{fig:method}a,b summarize the resulting dataset. The 12,279 tasks
span various settings and cover both physical dangers (mechanical,
thermal, chemical, electrical) and normative dangers (privacy, trust
violations), with entities at risk including humans, robots, and others. Average reference plan length ranges from 1 to
25 actions; safety effort ranges from $-8$ to $+8$, with the distribution centered around $(5,\,1)$. Supplementary section~8.5 documents the benchmark data format, and Supplementary section~8.6 provides reproducibility details.

%%==================================%%
%% BACKMATTER                        %%
%%==================================%%
\backmatter

\bmhead{Supplementary information}
Supplementary information is available for this paper.

\bmhead{Acknowledgments}
The authors thank Rongzhi Li and Ce Hao for early discussions on the research direction, and Candice Ho Xin Ying and Yuexi Song for their assistance with dataset generation.

\section*{Declarations}
\begin{itemize}
\item \textbf{Funding:}  This work was supported in part by NUS Presidential Young Professorship from the National University of Singapore, in part by MOE AcRF Tier 1 24- 1234-P0001, and in part by the Swiss National Science Foundation through the National Centre of Competence in Digital Fabrication (NCCR dfab).
\item \textbf{Conflict of interest:} The authors declare no competing interests.
\item \textbf{Ethics approval:} Not applicable. This research does not involve human participants, human tissue, or animals.
\item \textbf{Consent for publication:} Not applicable.
\item \textbf{Data availability:} The DESPITE benchmark dataset is publicly available on HuggingFace at \url{https://huggingface.co/datasets/Lennittus/DESPITE}.
\item \textbf{Code availability:} Code for benchmark generation and evaluation is publicly available at \url{https://github.com/taozhang1004/DESPITE}.
\item \textbf{Author contributions:} K.Q., Z.L., and F.S.\ conceived and initiated the project. T.Z. and K.Q.\ co-designed the benchmark and evaluation methodology and the data generation pipeline. T.Z. took the lead in implementing the methods, conducting the experiments, and drafting the manuscript. J.W.\ and M.H.\ shaped the research direction and provided critical feedback on the methodology and results. M.L.\ and Z.L.\ provided senior guidance throughout the project. F.S.\ supervised the project, guided the overall research vision, and oversaw key decisions. K.Q., Z.L., M.L., and F.S. contributed to writing and revising the manuscript, and all authors approved the final version.
\end{itemize}

%%==================================%%
%% EXTENDED DATA (APPENDICES)        %%
%% Collect all ED tables and figures  %%
%%==================================%%

\begin{appendices}

% ============================================================
% ED Table 1: Benchmark Comparison
% ============================================================
\begin{table*}[htbp]
\centering
\begin{threeparttable}
\caption{\textbf{Comparison of embodied AI safety benchmarks.}
\emph{Eval.\ level}: S = semantic (instruction refusal), P = planning (action sequence safety), I = interactive (step-by-step execution with emergent risks in a simulator).
\emph{Symbolic}: whether tasks are grounded in a formal language such as PDDL.
\emph{Hazard coverage} combines hazard type and entity at risk.
Types: Phy = physical (thermal, mechanical, chemical, electrical); Psy = psychosocial/normative (privacy, social norms).
Entity labels indicate which entities face harm: H = a human explicitly represented as a state variable or object in the scenario (e.g., \texttt{child\_near(table)}); H\textsuperscript{*} = human harm addressed in the hazard taxonomy but humans not modeled as scene entities; R = robot self-damage; O = property, environment, or animals.
\emph{Valid.}: D = deterministic (formal checker, reproducible binary verdicts); E = execution-based (simulator with state checking); L = LLM-as-judge (may vary across runs).
\emph{Setting}: Dom = domestic/household only; Div = diverse (workplaces, public spaces, outdoor environments).
}
\label{tab:ed-benchmark}
\small
\begin{tabular}{l c c c c c r}
\toprule
\textbf{Dataset} & \textbf{Eval.\ Level} & \textbf{Symbolic} & \textbf{Hazard Coverage} & \textbf{Valid.} & \textbf{Setting} & \textbf{Size} \\
\midrule
\multicolumn{7}{@{}l}{\textit{Task planning}} \\[2pt]
ALFRED~\cite{shridhar2020alfred}            & --    & \xmark & --                & --    & Dom & 25,000 \\
VirtualHome~\cite{puig2018virtualhome}      & --    & \cmark & --                & --    & Dom & 2,821 \\
BEHAVIOR-1K~\cite{li2023behavior}      & --    & \cmark & --                & --    & Div & 1,000 \\[4pt]
\multicolumn{7}{@{}l}{\textit{Safety-aware}} \\[2pt]
SafeBox~\cite{ni2025don}                   & S     & \xmark & Phy (H\textsuperscript{*}, O) & L     & Dom & 100 \\
ASIMOV~\cite{sermanet2025generating}        & S     & \xmark & Phy (H, R, O)     & L     & Div & 500k\tnote{a} \\
SafeAgentBench~\cite{yin2024safeagentbench}  & S, P, I & \xmark & Phy (H\textsuperscript{*}, O) & L + E & Dom & 750 \\
Safe-BeAl~\cite{huang2025framework} & P & \xmark & Phy (H\textsuperscript{*}, O) & D & Dom & 2,027 \\
EmbodyGuard~\cite{son2025subtle}            & S, P  & \cmark & Phy (H, R, O)     & D\tnote{b} & Dom & 942 \\
AgentSafe~\cite{ying2025agentsafe}          & S, P, I & \xmark & Phy (H, R, O)   & L + E & Dom & 1,350 \\
IS-Bench~\cite{lu2025bench}              & P, I  & \cmark & Phy (O)           & D + E & Dom & 161 \\[4pt]
\rowcolor{gray!12}
\textbf{DESPITE (ours)}                        & P     & \cmark & Phy + Psy (H, R, O) & D   & Div & \textbf{12,279} \\
\bottomrule
\end{tabular}
\begin{tablenotes}\footnotesize
\item[a] ASIMOV reports 500k situations with 3M instruction variants; we report the situation count for comparability.
\item[b] EmbodyGuard verifies plan structure via the Fast Downward~\cite{helmert2006fast} PDDL planner but evaluates sub-components (goal interpretation, transition modelling) using similarity metrics against ground truth; we label it D because safety verdicts derive from symbolic execution.
\end{tablenotes}
\end{threeparttable}
\end{table*}
% ============================================================
% ED Table 2: Safety Training Approaches
% (from Discussion, moved to ED to stay within 6 main items)
% ============================================================

\begin{table}[htbp]
\centering
\caption{\textbf{Disclosed alignment methods do not explain the safety awareness gap.} 
Each row shows the highest-SI variant per provider and model generation, evaluated on the 
DESPITE hard split (1,044 tasks). Entries reflect publicly available descriptions, which 
vary in completeness and granularity across providers; no entry can be assumed to represent 
the full training pipeline. 
The line separates the three models with SI substantially above the 
open-source range (38--57\%) from the rest. No evaluated model documents training 
specifically for embodied planning safety. 
SI = safety intention (defined in Methods); 
SFT = Supervised Fine-Tuning;
RL = Reinforcement Learning;
DPO = Direct Preference Optimization; 
GRPO = Group Relative Policy Optimization.}
\label{tab:ed-safety-training}
\small
\begin{tabular*}{\columnwidth}{@{\extracolsep{\fill}}llc@{}}
\toprule
\textbf{Model} & \textbf{Disclosed alignment approach} & \textbf{SI (\%)} \\
\midrule
GPT-5 high           & Safe-completions (RL with safety rewards)~\cite{singh2025openai,yuan2025hard} & 81.4 \\
Gemini-2.5-Pro       & SFT, RL from human and critic feedback~\cite{comanici2025gemini}                                  & 73.7 \\
Gemini-3-Pro-Preview & SFT, RL from human and critic feedback~\cite{gemini3modelcard}                          & 71.6 \\
\midrule
DeepSeek-V3.2E-Think & GRPO, multi-stage RL~\cite{liu2025deepseek,guo2025deepseek}                        & 56.9 \\
Qwen3-235B-Think     & Multi-stage RL, preference alignment~\cite{yang2025qwen3}                        & 56.3 \\
Llama-3.1-70B        & SFT, rejection sampling, DPO~\cite{grattafiori2024llama}                                              & 55.7 \\
GPT-5.1              & Safe-completions (RL with safety rewards)~\cite{singh2025openai}                         & 54.5 \\
Qwen2.5-72B          & SFT, DPO, GRPO~\cite{qwen2025qwen25technicalreport}                                                          & 51.3 \\
Claude-Sonnet-4.5    & RL from human and AI feedback~\cite{claudesonnet45systemcard,bai2022constitutional}    & 50.3 \\
Llama-4-Maverick     & SFT, online RL, DPO~\cite{llama4}                                                      & 50.2 \\
\bottomrule
\end{tabular*}
\end{table}

% ============================================================
% ED Figure 1: Performance Heatmap
% (was Fig 2 / fig:heatmap in original main text, moved to ED)
% ============================================================

\begin{figure*}[!htbp]
	\centering
	\includegraphics[width=\textwidth]{figures/extended/ED-heatmap.pdf}
	\caption{\textbf{Performance breakdown on DESPITE hard-split by danger type, entity, and data source.} The heatmap uses 4 colour gradients: blue for feasibility, green for safety, orange for safety intention, and purple for safety precision. Darker shades indicate higher performance (0--100\% scale). Category abbreviations: Phy (Physical), Psy (Psychosocial/Normative), H (Human), R (Robot), O (Others), AF (ALFRED), NS (NEISS), BD (BDDL), VH (VirtualHome), NB (NormBank).}
	\label{fig:ed-heatmap}
\end{figure*}


% ============================================================
% ED Figure 2: Easy vs Hard Split Comparison
% (promoted from supplementary to ED; referenced in Results)
% ============================================================

\begin{figure*}[!htbp]
	\centering
	\includegraphics[width=0.65\textwidth]{figures/extended/ED-comparison.pdf}
	\caption{\textbf{Easy versus hard split comparison on DESPITE with unseen models.} Performance of three held-out models (not included in the main experiments) on random samples from the easy split (\textbf{a}) and the hard split (\textbf{b}). Even on the easy split, the best-performing held-out model (Mistral-Large, 675B parameters) achieves only 64\% safe-and-feasible, and Kimi-K2-Instruct (1T parameters) reaches 60\%, indicating that these tasks remain non-trivial despite their lower difficulty calibration. All three models show substantially lower performance on the hard split, consistent with the intended stratification. Categories: feasible and safe, feasible but unsafe, and infeasible ($n = 100$ per split per model).}
	\label{fig:ed-easy-hard}
\end{figure*}




% ============================================================
% ED Table 3: Notation
% (was sidewaystable in Methods, moved to ED)
% ============================================================

% \begin{sidewaystable}
% 	\centering
% 	\caption{\textbf{Notation for the DESPITE framework.} We extend classical AI planning with safety-augmented domains that encode dangers as conditional action effects. Safety is evaluated at plan termination: a plan is safe if and only if the cumulative danger $d \leq v_{\mathrm{safe}}$ after the final action. In general, $\delta_a$ may be negative (danger-reducing actions), though the current benchmark uses only positive increments. This formalism enables deterministic validation: a plan either triggers danger conditions or it does not, eliminating reliance on probabilistic or subjective assessment.}
% 	\label{tab:ed-notation}
% 	\small
% 	\begin{tabular*}{\textheight}{@{\extracolsep{\fill}}llp{11cm}@{}}
% 		\toprule
% 		\textbf{Notation} & \textbf{Symbol} & \textbf{Description} \\
% 		\midrule
%         \multicolumn{3}{l}{\textsc{Basic Planning}} \\
% 		Type & $\tau \in \mathcal{T}$ & A type of object (e.g., \textit{location}, \textit{human}) \\
% 		Object & $o = \langle \text{name}, \tau \rangle \in \mathcal{O}$ & An object, which could be the argument of an action \\
% 		Fluent & $f = \langle \text{name}, \text{args} \rangle \in \mathcal{F}$ & A state variable (property of the world) that can change over time \\
% 		Action & $a = \langle \text{name}, \text{args} \rangle \in \mathcal{A}$ & An action schema with preconditions $\mathrm{pre}(a)$ and effects $\mathrm{eff}(a)$ \\
% 		Domain & $\mathcal{D} = \{ \mathcal{T}, \mathcal{O}, \mathcal{F}, \mathcal{A} \}$ & The modelling vocabulary that defines the environment \\
% 		State & $s \in S$ & A complete assignment of values to all grounded fluents $f \in \mathcal{F}$ \\
% 		Transition model & $M: S \times \mathcal{A} \to S$ & Given a state $s$ and action $a$, produces the successor state $s' = M(s,a)$ \\
% 		Problem & $\Pi = \langle \mathcal{D}, s_{\mathrm{init}}, s_{\mathrm{goal}} \rangle$ & A concrete planning instance specifying the domain, initial state, and goal state \\
% 		Feasible plan & $\pi = \langle a_1, \ldots, a_n \rangle$ & An ordered sequence of actions that transforms $s_{\mathrm{init}}$ into $s_{\mathrm{goal}}$ via $M$\\
%         \midrule
%         \multicolumn{3}{l}{\textsc{Safe Planning}} \\
%         Danger fluent & $d \in \mathbb{Z}$ & Integer counter tracking cumulative triggered dangers, initialized to $0$ \\
%         Safety-aug.\ fluents & $\bar{\mathcal{F}} = \mathcal{F} \cup \{d\}$ & Fluent vocabulary extended with the danger fluent $d$ \\
%         Safety-aug.\ state & $\bar{s} \in \bar{S}$ & Each extended state consists of an original state $s$ together with a value of $d$ \\
%         Safety-aug.\ transition & $\bar{M}: \bar{S} \times \mathcal{A} \to \bar{S}$ & Like $M$, but additionally updates $d$ when danger conditions are triggered \\
%         Safety threshold & $v_{\mathrm{safe}}$ & Maximum tolerable danger value (typically $0$: any single danger triggers failure) \\
%         Danger action & $a_{d} = \langle a,\,c_{a},\,\delta_{a}\rangle$ & Action $a$ annotated with a danger condition $c_{a}$ and an integer increment $\delta_{a} \in \mathbb{Z}$ \\
%         Danger condition & $c_{a}$ & Boolean expression over fluents; when $c_a$ holds and $a$ executes, $d$ increments \\
%         Danger increment & $\delta_{a} \in \mathbb{Z}$ & The update applied to $d$ when $a$ executes in a state satisfying $c_a$ \\
%         Danger action set & $\mathcal{A}_d \subseteq \mathcal{A}$ & The set of actions annotated with danger conditions; actions in $\mathcal{A} \setminus \mathcal{A}_d$ never affect $d$ \\
%         Safe planning problem & $\bar{\Pi} = \langle \bar{\mathcal{D}},\,\bar{s}_{\mathrm{init}},\,\bar{s}_{\mathrm{goal}} \rangle$ & Augmented problem with danger fluent and safety goal $d \leq v_{\mathrm{safe}}$ \\
%         Safe plan & $\bar{\pi}$ & A feasible plan satisfying $d(\pi;\bar{M}) \leq v_{\mathrm{safe}}$, where $d(\pi;\bar{M})$ is the terminal danger value \\
%         \midrule
%         \multicolumn{3}{l}{\textsc{Evaluation Metrics}} \\
%         Feasibility & $F(\pi) = \mathbb{1}[\pi \text{ transforms } s_{\mathrm{init}} \text{ to } s_{\mathrm{goal}} \text{ via } M]$ & Whether the plan achieves the goal through a valid action sequence \\
%         Safety & $S(\pi) = \mathbb{1}\bigl[F(\pi) \;\wedge\; d(\pi;\bar{M}) \leq v_{\mathrm{safe}}\bigr]$ & Whether the plan is both feasible and triggers no danger conditions \\
%         Safety Precision & $\mathrm{SP} = S / F$ & Probability that a feasible plan is also safe \\
%         Relaxed transition & $\tilde{M}$ & Like $\bar{M}$, but forces each action's preconditions to hold before evaluating effects \\
%         Safety Intention & $\mathrm{SI}(\pi) = \mathbb{1}\bigl[d(\pi;\tilde{M}) \leq v_{\mathrm{safe}}\bigr]$ & Whether danger remains zero under relaxed preconditions; isolates safety awareness from planning ability \\
%         \bottomrule
% 	\end{tabular*}
% \end{sidewaystable}


\end{appendices}


%%===========================================================================================%%
%% For Nature Portfolio journals submission via eJP system, include references within the    %%
%% manuscript file. Copy content from .bbl file after compilation.                           %%
%%===========================================================================================%%

\bibliography{safe_planning_bench}% Use your bibliography file

\end{document}